\begin{table}[tbp]
\centering
\caption{Behavioural responses and perceptions reported in interviews.}\label{tab:ch6-themes}
\fontsize{9.5}{12}\selectfont
\renewcommand{\arraystretch}{1.13}
\begin{tabularx}{\textwidth}{@{}>{\raggedright\arraybackslash}Xr@{}}
\toprule
Theme & Mentions\\
\midrule
\multicolumn{2}{@{}l}{\textsf{Communication}}\\
Move discussion from private channels into public records & 13\\
Use thumbs-up reactions for appreciation or acknowledgement & 9\\
Seek collaboration with contributors with higher OpenRank & 3\\
Use scores as supporting information when assigning issues & 2\\
Edit or consolidate comments instead of adding actions & 1\\
\multicolumn{2}{@{}l}{\textsf{Development}}\\
Split large changes into smaller, independent pull requests & 8\\
Use pull requests in place of direct commits & 5\\
Wait for others to review or merge changes & 3\\
Submit small fixes or dependency updates & 3\\
\multicolumn{2}{@{}l}{\textsf{Perception}}\\
Endorse contribution evaluation through collaboration networks & 11\\
Use the leaderboard for self-assessment and collaboration & 8\\
Perceive algorithmic complexity as raising the cost of gaming & 4\\
Suggest incorporating the value of contribution content & 2\\
View content-agnostic evaluation as encouraging newcomers & 1\\
\bottomrule
\end{tabularx}
\tablenote{Fourteen themes and their frequencies translated and consolidated from \citet[Table 6.7]{zhao2025thesis} and \citet[Table 6.7]{xia2024thesis}, corresponding to the same interview study reported in \citet[Section 4.2]{zhao2024}. The sample comprises 18 interviewees; themes overlap. Counts are mentions within this selected sample, not prevalence estimates for project populations. The theses and paper are overlapping accounts, not independent replications.}
\end{table}
